\documentclass[11pt]{article}
\usepackage[margin=1in]{geometry}
\usepackage{amsmath,amssymb,amsthm,hyperref,enumitem}
\newtheorem{theorem}{Theorem}
\newtheorem{lemma}{Lemma}
\newcommand{\E}{\mathbb E}
\newcommand{\R}{\mathbb R}
\title{Convexity of hypergraph polynomials from exponential minima}
\author{Principia Math In Progress}
\date{21 September 2026}
\begin{document}
\maketitle
\noindent\textbf{Task 1: MathDB 347913.}
Let $E=\binom{[n]}L$, $m=|E|$, and
\begin{equation}\label{definition}
 f_d(x)=\sum_{(e_1,\ldots,e_d)\in E^d}
         \prod_{k=1}^d\frac{x_{e_k}}{|e_1\cup\cdots\cup e_k|}.
\end{equation}
Conjecture 1 of Brosch, Laurent and Steenkamp \cite{BLS} asks whether
$f_d$ is convex on $\Delta_m=\{x\geq0:\sum_e x_e=1\}$ for every
$n,L,d\geq2$ (with $L\leq n$, so the simplex is nonempty).
We prove this, and in fact prove strict convexity on the whole
nonnegative orthant for any family of distinct nonempty hyperedges.

\begin{theorem}\label{main}
Let $V$ be a finite set and $E$ a nonempty family of distinct nonempty
subsets of $V$. Define $f_d$ by \eqref{definition}, with $d\geq2$.
Let $(T_v)_{v\in V}$ be independent exponential random variables
of rate one, and set $Y_e=\min_{v\in e}T_v$. Then
\begin{equation}\label{moment}
 \boxed{\displaystyle f_d(x)=\frac1{d!}\E\left[
                      \left(\sum_{e\in E}x_eY_e\right)^d\right]}
 \qquad (x\in\R^E).
\end{equation}
Consequently $f_d$ is strictly convex on $\R_+^E$. Its Hessian is
positive definite at every nonzero $x\in\R_+^E$; for $d=2$ it is
positive definite at the origin as well. When $E=\binom{[n]}L$,
the uniform vector is the unique minimizer on $\Delta_m$.
\end{theorem}

\begin{lemma}\label{mixed}
For any ordered tuple $(e_1,\ldots,e_d)$ of nonempty edges, including
tuples with repeated edges,
\begin{equation}\label{joint}
 \E\prod_{i=1}^dY_{e_i}
 =\sum_{\pi\in S_d}\prod_{k=1}^d
       \frac1{|e_{\pi(1)}\cup\cdots\cup e_{\pi(k)}|}.
\end{equation}
\end{lemma}
\begin{proof}
By the tail-integral formula and Tonelli's theorem, the left side is
\[
 \int_{[0,\infty)^d}\mathbb P(Y_{e_i}>s_i\text{ for all }i)\,ds
 =\int_{[0,\infty)^d}
    \exp\left(-\sum_{v\in V}\max_{i:v\in e_i}s_i\right)\,ds,
\]
where an empty maximum is zero. This probability formula follows
because $T_v$ must exceed every $s_i$ for which $v\in e_i$.
Partition the integration region, up to measure-zero ties, into the
$d!$ chambers $s_{\pi(1)}\geq\cdots\geq s_{\pi(d)}\geq0$.
In one chamber put $u_k=s_{\pi(k)}-s_{\pi(k+1)}$ and
$s_{\pi(d+1)}=0$. The Jacobian is one, and
\[
 \sum_{v\in V}\max_{i:v\in e_i}s_i
 =\sum_{k=1}^d |e_{\pi(1)}\cup\cdots\cup e_{\pi(k)}|\,u_k.
\]
Indeed, a vertex first covered at position $j$ contributes
$s_{\pi(j)}=u_j+\cdots+u_d$.
Integrating the $u_k$ separately gives the product in \eqref{joint}.
Permutations here order labeled positions, even when some $e_i$
coincide; all $d!$ chambers must still be included.
\end{proof}

\begin{proof}[Proof of Theorem \ref{main}]
All required moments are finite: $0<Y_e\leq\sum_{v\in V}T_v$
almost surely, and a finite sum of exponentials has all integer
moments. Expand $(\sum_e x_eY_e)^d$ as a sum over ordered edge
tuples and use Lemma \ref{mixed}. For each fixed $\pi\in S_d$,
the relabeling $(e_1,\ldots,e_d)\mapsto(e_{\pi(1)},\ldots,e_{\pi(d)})$
is a bijection of $E^d$ and leaves the monomial unchanged.
Thus every one of the $d!$ summands equals $f_d(x)$, proving
\eqref{moment}. This is a polynomial identity, including signed $x$.

Differentiating this finite polynomial identity gives
\begin{equation}\label{hessian}
 \nabla^2 f_d(x)=\frac1{(d-2)!}
            \E\big[(x\mathbin{\cdot}Y)^{d-2}YY^{\mathsf T}\big].
\end{equation}
For $d=2$ the power is identically one. For $x\geq0$ the weight
is nonnegative, so the right side is positive semidefinite.
This already proves the conjectured convexity, including the
boundary, by restriction to line segments.

For strictness we prove that the functions
$t\mapsto\min_{v\in e}t_v$, $e\in E$, are linearly independent
on $(0,\infty)^V$. If a linear combination with coefficients $a_e$
vanishes there, continuity extends the identity to $[0,\infty)^V$.
At the indicator vector $t=\mathbf1_S$, for $S\subseteq V$, it reads
\[
                 \sum_{e\in E:\ e\subseteq S}a_e=0.
\]
Extend $a_e$ by zero to subsets outside $E$. Induction on $|S|$
(equivalently, Boolean-lattice M\"obius inversion) shows every
$a_e=0$. Since the exponential joint density is positive throughout
$(0,\infty)^V$, continuity implies that for every nonzero $a$,
$\mathbb P(a\mathbin{\cdot}Y\neq0)>0$.
For nonzero $x\geq0$, we have $x\mathbin{\cdot}Y>0$ almost surely.
Equation \eqref{hessian} is therefore positive definite.
Alternatively, strict convexity on the entire orthant, including
segments from the origin, follows directly from the strict convexity
of $t\mapsto t^d$ on $[0,\infty)$: for distinct $x,z\geq0$,
$x\mathbin{\cdot}Y\neq z\mathbin{\cdot}Y$ with positive probability.

Finally take $E=\binom{[n]}L$. Permuting vertices preserves
\eqref{definition} and acts transitively on the edge set.
Averaging any $x\in\Delta_m$ over all vertex permutations gives
$x^*=(1/m,\ldots,1/m)$. Convexity and invariance imply
$f_d(x^*)\leq f_d(x)$, and strict convexity gives uniqueness.
This last symmetry argument is the one already used in \cite{BLS}.
\end{proof}

\paragraph{The coefficient matrices in the source.}
Write $f_d(x)=\sum_{|\alpha|=d}b_\alpha x^\alpha$ and
$\alpha!=\prod_e\alpha_e!$. Formula \eqref{moment} implies
$b_\alpha=\E[Y^\alpha]/\alpha!$. Thus the matrices in
\cite[Lemmas 4.1 and 4.3]{BLS} satisfy
\begin{equation}\label{gram}
 \nabla^2f_d(x)=\sum_{|\gamma|=d-2}x^\gamma Q_\gamma,
 \qquad Q_\gamma=\frac1{\gamma!}\E[Y^\gamma YY^{\mathsf T}].
\end{equation}
They are positive definite, since $Y^\gamma>0$ almost surely and
the independence argument above applies. This proves the stronger
coefficientwise positivity question posed in the source as well.
For a single edge of size $L$, \eqref{moment} gives
$f_d(x)=x^d/L^d$, confirming the factor $d!$.

\paragraph{Independent finite verification.}
The supplied standard-library Python verifier uses exact rational
arithmetic. One algorithm evaluates the original ordered-edge sum
by merging prefixes with the same union. A second integrates over
the $|V|!$ chambers of ordered \emph{vertex} lifetimes. To describe
the second algorithm, put $N=|V|$, let
$T_{\sigma(1)}\leq\cdots\leq T_{\sigma(N)}$, and define
$S_j=\{\sigma(j),\ldots,\sigma(N)\}$,
$h_j=\sum_{e\subseteq S_j}x_e$, $r_j=N-j+1$. Then
\[
 f_d(x)=\frac1{N!}\sum_{\sigma\in S_N}
 [z^d]\prod_{j=1}^N\frac1{1-z h_j/r_j}.
\]
Indeed, on this chamber, successive spacings $u_j\geq0$ have
density $\exp(-\sum_jr_ju_j)$ and $x\cdot Y=\sum_jh_ju_j$.
Expanding its $d$th power and integrating uses
$\int_0^\infty u^k e^{-ru}\,du=k!/r^{k+1}$;
the product of the $r_j$ is $N!$.

The two algorithms agree in 585 exact value comparisons through
degree 12 and in 596 symbolic coefficients, including repeated
edges, mixed edge sizes, signed inputs, and simplex vertices.
All 84 sampled Hessian coefficient matrices have exact positive
elimination pivots. Separate checks verify the single-edge
normalization and the Boolean inversion used for strictness.
These finite checks supplement the unrestricted proof above.

\paragraph{Scope and attribution.}
The source proves convexity for $d=2$ and for $(d,L)=(3,2)$;
Polak \cite{Polak} proves $L=2$, $d\leq9$.
The present contribution is the moment representation and its
application for arbitrary $d$ and nonempty edge families.
Tail integration, exponential survival probabilities, permutation
chambers, and Gram positivity are classical tools.
Both the complete current target preprint and its journal final
were read. The bounded prior-claim audit and its limits are recorded
separately. Review is by the authoring agent only. The minimum
corollary and matrix positivity are included in this one result;
they receive no separate campaign credit.

\begin{enumerate}[label=\textbf{Task \arabic*:},start=2,leftmargin=*]
\item Find useful explicit lower bounds for the smallest eigenvalue
of $Q_\gamma$ in terms of the hypergraph and $\gamma$.
\item Determine sharp quantitative gaps between $f_d(x)$ and its
uniform minimum for complete uniform hypergraphs.
\item Use the moment representation to study generating functions
in $d$ and their relation to queue stability and occupancy.
\item For unequal vertex service rates, characterize the minimizing
edge distribution and the symmetries that force its coordinates.
\end{enumerate}

\small
\begin{thebibliography}{9}
\raggedright
\bibitem{BLS} D. Brosch, M. Laurent and A. Steenkamp,
Optimizing hypergraph-based polynomials modeling job-occupancy in
queuing with redundancy scheduling,
\emph{SIAM Journal on Optimization} \textbf{31}(3) (2021), 2227--2254.
\href{https://doi.org/10.1137/20M1369592}{doi:10.1137/20M1369592}.
Conjecture 1, journal p.~2242; source matrices on pp.~2243--2244.
\href{https://arxiv.org/abs/2009.04510}{arXiv:2009.04510v1},
Conjecture 1 on preprint p.~20.
\bibitem{Polak} S. C. Polak, Symmetry reduction to optimize a
graph-based polynomial from queueing theory,
\emph{SIAM Journal on Applied Algebra and Geometry}
\textbf{6} (2022), 243--266.
\href{https://doi.org/10.1137/21M1413298}{doi:10.1137/21M1413298}.
\href{https://arxiv.org/abs/2104.08264}{arXiv:2104.08264v2},
Theorem 1.1.
\end{thebibliography}
\end{document}
